\clearpage\thispagestyle{empty}
\pdfbookmark[0]{Quellen zu den XIII. Elementen der Ceres (3. Juni 1809)}{ceres1809xiii-sources}
\begin{center}\textbf{Quellen zu den XIII. Elementen der Ceres (3. Juni 1809)}\end{center}
\noindent Diese KI-Übertragung ins heutige Deutsch folgt der Mitteilung zur Ceres in \href{https://gdz.sub.uni-goettingen.de/id/PPN236007467}{Gauß’ \textit{Werken}, Bd.~VI} (Göttingen, 1874), S.~312--314. Der Erstabdruck erschien unter \textit{Göttingen} in den \textit{Göttingischen gelehrten Anzeigen}, Stück 87, 3. Juni 1809, S.~\href{https://api.digitale-sammlungen.de/iiif/image/v2/bsb10538675_00859/full/1600,/0/default.jpg}{857}, \href{https://api.digitale-sammlungen.de/iiif/image/v2/bsb10538675_00860/full/1600,/0/default.jpg}{858}, \href{https://api.digitale-sammlungen.de/iiif/image/v2/bsb10538675_00861/full/1600,/0/default.jpg}{859}, \href{https://api.digitale-sammlungen.de/iiif/image/v2/bsb10538675_00862/full/1600,/0/default.jpg}{860}. Das digitalisierte Exemplar stammt aus der Bayerischen Staatsbibliothek.
\par\smallskip\noindent Siehe auch Bernhard August von Lindenaus \href{https://archive.org/download/bub_gb_caMAAAAAMAAJ/page/n518.jpg}{\textit{Fortgesetzte Nachrichten über die neuen Hauptplaneten}}, \textit{Monatliche Correspondenz}, Bd.~XIX, Mai 1809, S.~504--515, im von Google digitalisierten Exemplar der New York Public Library. Der Parallelbericht über Ceres steht auf S.~504--507. In der \href{https://archive.org/download/bub_gb_caMAAAAAMAAJ/page/n525.jpg}{Ceres-Tafel} bezeichnet die letzte Spalte den Logarithmus des Abstands, in der \href{https://archive.org/download/bub_gb_caMAAAAAMAAJ/page/n526.jpg}{Juno-Tafel} den Abstand von der Erde. Die eckige Klammer folgt dem Wortlaut der Werkausgabe.
\par\smallskip\noindent KI-Übertragung und KI-Satz.
\clearpage
